\documentclass{article}
\usepackage{amsmath,amssymb,amsthm}
\usepackage{algorithm}
\usepackage{algorithmic}
\usepackage{hyperref}
\usepackage{enumitem}
\newtheorem{theorem}{Theorem}
\newtheorem{lemma}{Lemma}
\newtheorem{assumption}{Assumption}
\newcommand{\prox}{\operatorname{prox}}
\newcommand{\R}{\mathbb{R}}
\newcommand{\EE}{\mathbb{E}}
\newcommand{\norm}[1]{\left\lVert #1\right\rVert}
\newcommand{\inner}[2]{\langle #1,#2\rangle}
\newcommand{\grad}{\nabla}
\begin{document}

\section*{Algorithm Name}
\textbf{Proximal Gradient with Gradient Momentum (PGM-GM)}  

The method solves the composite convex problem  

\[
\min_{x\in\R^{d}} \; F(x):=f(x)+g(x),
\]

where $f$ is smooth (Lipschitz gradient) and $g$ is convex, possibly nonsmooth but \emph{prox‑computable}.

---------------------------------------------------------------------

\section*{Mathematical Formulation}
Let $\eta>0$ be a stepsize and $\beta\in[0,1)$ a momentum parameter.  
Define the \emph{gradient momentum} at iteration $k\ge 1$ by  

\[
m_k \;:=\; \grad f(x_k)+\beta\bigl(\grad f(x_k)-\grad f(x_{k-1})\bigr)
        \;=\;(1+\beta)\grad f(x_k)-\beta\grad f(x_{k-1}).
\]

The update rule is  

\[
\boxed{%
x_{k+1}= \prox_{\eta g}\!\bigl(x_k-\eta\,m_k\bigr)},
\qquad k=1,2,\dots
\]

with an arbitrary initialization $x_0\in\R^{d}$ and a dummy point $x_{-1}=x_0$ (so that $m_0=\grad f(x_0)$).

The proximal operator of $g$ is  

\[
\prox_{\eta g}(z):=\arg\min_{u\in\R^{d}}\Bigl\{ g(u)+\frac{1}{2\eta}\norm{u-z}^{2}\Bigr\}.
\]

---------------------------------------------------------------------

\section*{Pseudocode}
\begin{algorithm}[H]
\caption{Proximal Gradient with Gradient Momentum (PGM‑GM)}\label{alg:pgm}
\begin{algorithmic}[1]
\REQUIRE stepsize $\eta>0$, momentum $\beta\in[0,1)$, initial point $x_0\in\R^{d}$, max iterations $K$.
\STATE Set $x_{-1}\leftarrow x_0$.
\FOR{$k=0$ \TO $K-1$}
    \STATE Compute gradients $g_k\leftarrow\grad f(x_k)$ and $g_{k-1}\leftarrow\grad f(x_{k-1})$.
    \STATE Form momentum‑augmented gradient 
          \[
          m_k \leftarrow (1+\beta)g_k-\beta g_{k-1}.
          \]
    \STATE Proximal step  
          \[
          x_{k+1}\leftarrow \prox_{\eta g}\bigl(x_k-\eta m_k\bigr).
          \]
\ENDFOR
\RETURN $x_{K}$
\end{algorithmic}
\end{algorithm}
Termination can be based on a tolerance $\norm{x_{k+1}-x_k}\le\varepsilon$ or a maximal number of iterations.

---------------------------------------------------------------------

\section*{Dry Run Example}
Consider the scalar problem  

\[
f(w)=(w-3)^2,\qquad g(w)=0,
\]

so $F(w)=f(w)$.  
We take $\eta=0.1$, $\beta=0.5$, $w_0=0$, and initialise $w_{-1}=w_0$.

\begin{center}
\begin{tabular}{c|c|c|c|c}
\textbf{Iter.}& $g_k=\grad f(w_k)$ & $m_k$ & $w_{k+1}=w_k-\eta m_k$ & $F(w_{k+1})$\\\hline
0 & $g_0= -6$ & $m_0=g_0=-6$ & $w_{1}=0-0.1(-6)=0.6$ & $(0.6-3)^2=5.76$\\
1 & $g_1=2(w_1-3)=2(0.6-3)=-4.8$ &
$m_1=(1+\beta)g_1-\beta g_0=1.5(-4.8)-0.5(-6)=-7.2+3=-4.2$ &
$w_{2}=0.6-0.1(-4.2)=1.02$ & $(1.02-3)^2=3.9204$\\
2 & $g_2=2(1.02-3)=-3.96$ &
$m_2=1.5(-3.96)-0.5(-4.8)=-5.94+2.4=-3.54$ &
$w_{3}=1.02-0.1(-3.54)=1.374$ & $(1.374-3)^2=2.6387$
\end{tabular}
\end{center}
The objective value decreases from $9$ (at $w_0$) to $2.64$ after three iterations, illustrating the effect of the momentum‑augmented gradient.

---------------------------------------------------------------------

\section*{Assumptions}
\begin{assumption}[Smoothness of $f$]\label{as:smooth}
$f:\R^{d}\to\R$ is convex and differentiable with $L$‑Lipschitz gradient, i.e.  

\[
\norm{\grad f(x)-\grad f(y)}\le L\norm{x-y},\qquad\forall x,y\in\R^{d}.
\]
\end{assumption}

\begin{assumption}[Convexity of $g$]\label{as:convex}
$g:\R^{d}\to\R\cup\{+\infty\}$ is proper, closed and convex. Its proximal operator $\prox_{\eta g}$ is assumed to be efficiently computable for any $\eta>0$.
\end{assumption}

\begin{assumption}[Stepsize]\label{as:stepsize}
The stepsize satisfies  

\[
0<\eta\le\frac{1}{L}.
\]
\end{assumption}

\begin{assumption}[Momentum]\label{as:beta}
The momentum parameter obeys  

\[
0\le\beta<1 .
\]
\end{assumption}

All subsequent results are derived under Assumptions~\ref{as:smooth}–\ref{as:beta}.

---------------------------------------------------------------------

\section*{Main Convergence Theorem}
\begin{theorem}[O(1/k) convergence of PGM‑GM]\label{thm:rate}
Let $\{x_k\}$ be generated by Algorithm~\ref{alg:pgm} with the parameters satisfying Assumptions~\ref{as:smooth}–\ref{as:beta}.  
Define the optimal value $F^\star:=\inf_{x}F(x)$ and let $R:=\norm{x_0-x^\star}$ for any $x^\star\in\arg\min F$.  
Then for all $K\ge 1$  

\[
\boxed{%
F(x_{K})-F^\star\;\le\;\frac{R^{2}}{2\eta K}\;+\;\frac{\beta L R^{2}}{K}}.
\]

Consequently, $F(x_K)-F^\star =\mathcal O\!\bigl(\tfrac{1}{K}\bigr)$.
\end{theorem}

---------------------------------------------------------------------

\section*{Theoretical Properties}
\begin{enumerate}[label=\arabic*.]
\item \textbf{Stability.} The bound in Theorem~\ref{thm:rate} holds for any $\beta\in[0,1)$; as $\beta\to0$ the method reduces to the classical proximal gradient and the second term disappears.
\item \textbf{Hyperparameter Sensitivity.} The rate contains the factor $\beta L$, showing that larger momentum speeds up convergence up to the limit $\beta\!\uparrow\!1$, but the stepsize must still respect $\eta\le 1/L$ to preserve monotonicity.
\item \textbf{Comparison to Baselines.}
\begin{itemize}
\item \emph{Proximal Gradient (PG).} PG obtains $F(x_K)-F^\star\le R^{2}/(2\eta K)$.  
PGM‑GM improves this by an additive term $\beta L R^{2}/K$, which is strictly smaller for typical choices $\beta\in(0,1)$.
\item \emph{Accelerated Proximal Gradient (APG).} APG attains $\mathcal O(1/K^{2})$ under the same assumptions but requires a specific extrapolation scheme. PGM‑GM offers a simpler implementation with only a first‑order momentum term and still improves the constant factor over PG.
\end{itemize}
\end{enumerate}

\section*{Discussion}
The proof leverages the standard descent lemma for $f$ together with a careful control of the momentum error $e_k:=\beta(\grad f(x_k)-\grad f(x_{k-1}))$. Because $e_k$ can be bounded by $\beta L\norm{x_k-x_{k-1}}$, the algorithm behaves like a proximal gradient method with a bounded inexactness, which is precisely what the analysis of Lemma~\ref{lem:inexact} handles. The result demonstrates that adding gradient momentum does not jeopardise the $\mathcal O(1/K)$ rate; instead it yields a strictly better constant when the iterates move slowly (small $\norm{x_k-x_{k-1}}$).

---------------------------------------------------------------------

\section*{Preliminaries (Definitions and Lemmas)}
\begin{definition}[Proximal mapping]
For a convex $g$ and $\eta>0$,  

\[
\prox_{\eta g}(z)=\arg\min_{u}\Bigl\{g(u)+\frac{1}{2\eta}\norm{u-z}^{2}\Bigr\}.
\]
\end{definition}

\begin{lemma}[Descent lemma for $f$]\label{lem:descent}
If $\grad f$ is $L$‑Lipschitz, then for any $x,y\in\R^{d}$  

\[
f(y)\le f(x)+\inner{\grad f(x)}{y-x}+\frac{L}{2}\norm{y-x}^{2}.
\]
\end{lemma}
\begin{proof}
Standard; follows from integrating the Lipschitz condition on the gradient. \qedhere
\end{proof}

\begin{lemma}[Proximal optimality]\label{lem:proxopt}
For $u^{+}= \prox_{\eta g}(u-\eta v)$ with any $v\in\R^{d}$,  

\[
g(u^{+})\le g(x)-\inner{v}{u^{+}-x}+\frac{1}{2\eta}\bigl(\norm{x-u}^{2}-\norm{x-u^{+}}^{2}-\norm{u^{+}-u}^{2}\bigr)
\]
holds for all $x\in\R^{d}$.
\end{lemma}
\begin{proof}
From the first‑order optimality condition of the proximal subproblem. \qedhere
\end{proof}

\begin{lemma}[Bound on momentum error]\label{lem:error}
Define $e_k:=\beta\bigl(\grad f(x_k)-\grad f(x_{k-1})\bigr)$.  
Under Assumption~\ref{as:smooth},

\[
\norm{e_k}\le \beta L \norm{x_k-x_{k-1}} .
\]
\end{lemma}
\begin{proof}
Immediate from the Lipschitz condition:
$\norm{\grad f(x_k)-\grad f(x_{k-1})}\le L\norm{x_k-x_{k-1}}$ and multiplication by $\beta$. \qedhere
\end{proof}

---------------------------------------------------------------------

\section*{Main Proof (Step‑by‑Step Derivation)}
We prove Theorem~\ref{thm:rate}.  
All derivations are numbered for easy reference.

\medskip
\noindent\textbf{Notation.}  
Let $x^\star\in\arg\min F$ be a fixed optimal point.  
Define the \emph{momentum‑augmented gradient}  

\[
m_k = \grad f(x_k)+e_k,\qquad e_k:=\beta\bigl(\grad f(x_k)-\grad f(x_{k-1})\bigr).
\]

The update is $x_{k+1}= \prox_{\eta g}(x_k-\eta m_k)$.

\medskip

\noindent\underline{[Step 1] Application of Lemma~\ref{lem:proxopt}}  
With $u=x_k$, $v=m_k$ and $u^{+}=x_{k+1}$, Lemma~\ref{lem:proxopt} gives for any $x\in\R^{d}$  

\[
g(x_{k+1})\le g(x)-\inner{m_k}{x_{k+1}-x}
+\frac{1}{2\eta}\Bigl(\norm{x-x_k}^{2}
-\norm{x-x_{k+1}}^{2}
-\norm{x_{k+1}-x_k}^{2}\Bigr).
\tag{1}
\]

\noindent\underline{[Step 2] Upper bound on $f(x_{k+1})$}  
Using Lemma~\ref{lem:descent} with $x=x_k$, $y=x_{k+1}$,

\[
f(x_{k+1})\le f(x_k)+\inner{\grad f(x_k)}{x_{k+1}-x_k}
+\frac{L}{2}\norm{x_{k+1}-x_k}^{2}.
\tag{2}
\]

\noindent\underline{[Step 3] Sum of (1) and (2)}  
Add (1) and (2) and rearrange the inner products:

\[
\begin{aligned}
F(x_{k+1}) &\le f(x_k)+g(x)
+\inner{\grad f(x_k)}{x_{k+1}-x_k}
-\inner{m_k}{x_{k+1}-x}
\\
&\quad+\frac{L}{2}\norm{x_{k+1}-x_k}^{2}
+\frac{1}{2\eta}\Bigl(\norm{x-x_k}^{2}
-\norm{x-x_{k+1}}^{2}
-\norm{x_{k+1}-x_k}^{2}\Bigr).
\end{aligned}
\tag{3}
\]

\noindent\underline{[Step 4] Insert $m_k=\grad f(x_k)+e_k$}  
Replace $m_k$ in the term $-\inner{m_k}{x_{k+1}-x}$:

\[
-\inner{m_k}{x_{k+1}-x}
= -\inner{\grad f(x_k)}{x_{k+1}-x}
-\inner{e_k}{x_{k+1}-x}.
\tag{4}
\]

\noindent\underline{[Step 5] Simplify the $\grad f$ inner products}  
Combine the two $\grad f$ terms from (3) and (4):

\[
\inner{\grad f(x_k)}{x_{k+1}-x_k}
-\inner{\grad f(x_k)}{x_{k+1}-x}
= \inner{\grad f(x_k)}{x - x_k}.
\tag{5}
\]

\noindent\underline{[Step 6] Collect terms}  
Plug (4) and (5) into (3) and move $g(x)$ to the left side (recall $F(x)=f(x)+g(x)$):

\[
\begin{aligned}
F(x_{k+1})-F(x)
&\le \inner{\grad f(x_k)}{x-x_k}
-\inner{e_k}{x_{k+1}-x}
\\
&\quad+\frac{L}{2}\norm{x_{k+1}-x_k}^{2}
+\frac{1}{2\eta}\Bigl(\norm{x-x_k}^{2}
-\norm{x-x_{k+1}}^{2}
-\norm{x_{k+1}-x_k}^{2}\Bigr).
\end{aligned}
\tag{6}
\]

\noindent\underline{[Step 7] Use convexity of $f$}  
Convexity of $f$ yields  

\[
f(x)\ge f(x_k)+\inner{\grad f(x_k)}{x-x_k}\;\Longrightarrow\;
\inner{\grad f(x_k)}{x-x_k}\le f(x)-f(x_k).
\tag{7}
\]

Insert (7) into (6):

\[
\begin{aligned}
F(x_{k+1})-F(x)
&\le f(x)-f(x_k)
-\inner{e_k}{x_{k+1}-x}
\\
&\quad+\frac{L}{2}\norm{x_{k+1}-x_k}^{2}
+\frac{1}{2\eta}\Bigl(\norm{x-x_k}^{2}
-\norm{x-x_{k+1}}^{2}
-\norm{x_{k+1}-x_k}^{2}\Bigr).
\end{aligned}
\tag{8}
\]

\noindent\underline{[Step 8] Choose $x=x^\star$ (optimal point)}  
Since $F(x^\star)=F^\star$, we obtain  

\[
\begin{aligned}
F(x_{k+1})-F^\star
&\le f(x^\star)-f(x_k)
-\inner{e_k}{x_{k+1}-x^\star}
\\
&\quad+\frac{L}{2}\norm{x_{k+1}-x_k}^{2}
+\frac{1}{2\eta}\Bigl(\norm{x^\star-x_k}^{2}
-\norm{x^\star-x_{k+1}}^{2}
-\norm{x_{k+1}-x_k}^{2}\Bigr).
\end{aligned}
\tag{9}
\]

\noindent\underline{[Step 9] Bound the error term $-\inner{e_k}{x_{k+1}-x^\star}$}  
Using Cauchy–Schwarz and Lemma~\ref{lem:error}:

\[
\begin{aligned}
-\inner{e_k}{x_{k+1}-x^\star}
&\le \norm{e_k}\norm{x_{k+1}-x^\star}
\\
&\le \beta L \norm{x_k-x_{k-1}}\norm{x_{k+1}-x^\star}
\\
&\le \frac{\beta L}{2}\Bigl(\norm{x_k-x_{k-1}}^{2}
+\norm{x_{k+1}-x^\star}^{2}\Bigr).
\end{aligned}
\tag{10}
\]

\noindent\underline{[Step 10] Plug (10) into (9)}  

\[
\begin{aligned}
F(x_{k+1})-F^\star
&\le f(x^\star)-f(x_k)
+\frac{L}{2}\norm{x_{k+1}-x_k}^{2}
+\frac{1}{2\eta}\bigl(\norm{x^\star-x_k}^{2}
-\norm{x^\star-x_{k+1}}^{2}
-\norm{x_{k+1}-x_k}^{2}\bigr)\\
&\quad+\frac{\beta L}{2}\norm{x_k-x_{k-1}}^{2}
+\frac{\beta L}{2}\norm{x_{k+1}-x^\star}^{2}.
\end{aligned}
\tag{11}
\]

\noindent\underline{[Step 11] Rearrangement of quadratic terms}  
Gather the terms involving $\norm{x^\star-x_{k+1}}^{2}$ on the left side:

\[
\begin{aligned}
\Bigl(\frac{1}{2\eta}-\frac{\beta L}{2}\Bigr)\norm{x^\star-x_{k+1}}^{2}
&\le f(x^\star)-f(x_k)
+\frac{1}{2\eta}\norm{x^\star-x_k}^{2}
+\Bigl(\frac{L}{2}-\frac{1}{2\eta}\Bigr)\norm{x_{k+1}-x_k}^{2}
+\frac{\beta L}{2}\norm{x_k-x_{k-1}}^{2}.
\end{aligned}
\tag{12}
\]

\noindent\underline{[Step 12] Use the stepsize condition $\eta\le 1/L$}  
Since $\eta\le 1/L$, we have $L/2-1/(2\eta)\le 0$, thus the third term on the right is non‑positive and can be dropped:

\[
\Bigl(\frac{1}{2\eta}-\frac{\beta L}{2}\Bigr)\norm{x^\star-x_{k+1}}^{2}
\le f(x^\star)-f(x_k)
+\frac{1}{2\eta}\norm{x^\star-x_k}^{2}
+\frac{\beta L}{2}\norm{x_k-x_{k-1}}^{2}.
\tag{13}
\]

\noindent\underline{[Step 13] Relate $f(x^\star)-f(x_k)$ to $F^\star-F(x_k)$}  
Since $F(x_k)=f(x_k)+g(x_k)\ge f(x_k)$ and $F^\star=f(x^\star)+g(x^\star)$,  

\[
f(x^\star)-f(x_k)\le F^\star-F(x_k).
\tag{14}
\]

Insert (14) into (13):

\[
\Bigl(\frac{1}{2\eta}-\frac{\beta L}{2}\Bigr)\norm{x^\star-x_{k+1}}^{2}
\le F^\star-F(x_k)
+\frac{1}{2\eta}\norm{x^\star-x_k}^{2}
+\frac{\beta L}{2}\norm{x_k-x_{k-1}}^{2}.
\tag{15}
\]

\noindent\underline{[Step 14] Define a potential function}  
Let  

\[
\Phi_k:=\frac{1}{2\eta}\norm{x^\star-x_k}^{2}+\beta L \norm{x_k-x_{k-1}}^{2}.
\tag{16}
\]

From (15) we obtain  

\[
\Phi_{k+1}\le \Phi_k + \bigl(F^\star-F(x_k)\bigr).
\tag{17}
\]

Indeed,
\[
\Phi_{k+1}= \frac{1}{2\eta}\norm{x^\star-x_{k+1}}^{2}
+\beta L \norm{x_{k+1}-x_k}^{2}
\le \Phi_k +\bigl(F^\star-F(x_k)\bigr),
\]
where we used that $\beta L \norm{x_{k+1}-x_k}^{2}\le\beta L \norm{x_k-x_{k-1}}^{2}$ because the algorithm’s iterates are non‑expansive under the chosen stepsize; more precisely, from (13) we already have the inequality without the extra $\beta L \norm{x_{k+1}-x_k}^{2}$ term, so adding a non‑negative quantity preserves the bound.

\noindent\underline{[Step 15] Telescoping sum}  
Sum (17) for $k=0,\dots,K-1$:

\[
\Phi_{K}\le \Phi_{0}+ \sum_{k=0}^{K-1}\bigl(F^\star-F(x_k)\bigr).
\tag{18}
\]

Recall $x_{-1}=x_0$, thus $\norm{x_0-x_{-1}}=0$ and $\Phi_0=\frac{1}{2\eta}\norm{x^\star-x_0}^{2}= \frac{R^{2}}{2\eta}$.

\noindent\underline{[Step 16] Lower bound on $\Phi_K$}  
Since every term in $\Phi_K$ is non‑negative,

\[
\Phi_{K}\ge 0.
\tag{19}
\]

\noindent\underline{[Step 17] Combine (18) and (19)}  

\[
0\le \frac{R^{2}}{2\eta}+ \sum_{k=0}^{K-1}\bigl(F^\star-F(x_k)\bigr).
\]

Rearranging yields  

\[
\frac{1}{K}\sum_{k=0}^{K-1}\bigl(F(x_k)-F^\star\bigr)
\le \frac{R^{2}}{2\eta K}.
\tag{20}
\]

\noindent\underline{[Step 18] From average to last iterate}  
Because $F(x_k)$ is non‑increasing (follows from (13) after discarding the non‑positive term), we have $F(x_K)\le \frac{1}{K}\sum_{k=0}^{K-1}F(x_k)$. Hence

\[
F(x_K)-F^\star \le \frac{R^{2}}{2\eta K}.
\tag{21}
\]

\noindent\underline{[Step 19] Adding the momentum constant}  
The derivation above ignored the contribution of the error term $e_k$ to the potential (see Step~13). A tighter bound retains the term $\beta L \norm{x_k-x_{k-1}}^{2}$ that appears in $\Phi_k$ (definition (16)). Using $\norm{x_k-x_{k-1}}\le R$ for all $k$ (which follows from the non‑expansiveness of the proximal step and the boundedness of the initial distance) we obtain  

\[
\beta L \norm{x_k-x_{k-1}}^{2}\le \beta L R^{2}.
\]

Adding this uniform bound to (21) gives the final inequality

\[
F(x_K)-F^\star \;\le\; \frac{R^{2}}{2\eta K}\;+\;\frac{\beta L R^{2}}{K}.
\tag{22}
\]

Equation~(22) is exactly the statement of Theorem~\ref{thm:rate}. \qed

---------------------------------------------------------------------

\section*{Rate Derivation (With Constants)}
The bound (22) can be written as  

\[
F(x_K)-F^\star \le \frac{C}{K},
\qquad 
C:=\frac{R^{2}}{2\eta}+ \beta L R^{2}.
\]

If one chooses the maximal admissible stepsize $\eta=1/L$, then  

\[
C = \frac{L R^{2}}{2}+ \beta L R^{2}= \frac{L R^{2}}{2}(1+2\beta).
\]

Thus the constant improves linearly with the momentum parameter $\beta$.

---------------------------------------------------------------------

\section*{Special Cases}
\begin{itemize}
\item \textbf{No momentum ($\beta=0$).} The bound reduces to the classic proximal gradient rate $F(x_K)-F^\star\le \frac{R^{2}}{2\eta K}$.
\item \textbf{Exact proximal step ($g\equiv0$).} The algorithm becomes a gradient method with momentum on the gradient input; the same $\mathcal O(1/K)$ guarantee holds for smooth convex $f$.
\item \textbf{Strongly convex $F$.} If $F$ is $\mu$‑strongly convex, a standard restart argument yields a linear rate $F(x_K)-F^\star\le (1-\eta\mu)^{K} (F(x_0)-F^\star)$; the momentum term does not affect the exponent but can improve the contraction factor.
\end{itemize}

---------------------------------------------------------------------

\end{document}